\subsubsection{Building a Two-Qubit System}

In the previous sections, we have seen how to represent the state of a single qubit using a two-dimensional complex vector. Now, we want to extend this representation to a system of two qubits. To do this, we will use the concept of the tensor product seen in the previous section.

\highspace
We start from two independent qubits:
\begin{equation*}
    \ket{v_{A}} = a_{0}\ket{0} + a_{1}\ket{1} \quad \text{and} \quad \ket{v_{B}} = b_{0}\ket{0} + b_{1}\ket{1}
\end{equation*}
And our \textbf{goal is to find a state that encodes all joint possibilities} of the two qubits:
\begin{itemize}
    \item Both qubits are in state $\ket{0}$: $\ket{00}$
    \item The first qubit is in state $\ket{0}$ and the second in state $\ket{1}$: $\ket{01}$
    \item The first qubit is in state $\ket{1}$ and the second in state $\ket{0}$: $\ket{10}$
    \item Both qubits are in state $\ket{1}$: $\ket{11}$
\end{itemize}
To represent these joint possibilities, we will use the tensor product of the two qubits' states:
\begin{equation*}
    \ket{v_{A} v_{B}} = \ket{v_{A}} \otimes \ket{v_{B}}
\end{equation*}
Expanding this expression, we get:
\begin{equation}\label{eq:two-qubit-state}
    \begin{aligned}
        \ket{v_{A} v_{B}} &= (a_{0}\ket{0} + a_{1}\ket{1}) \otimes (b_{0}\ket{0} + b_{1}\ket{1}) \\
        &= a_{0}b_{0}\ket{00} + a_{0}b_{1}\ket{01} + a_{1}b_{0}\ket{10} + a_{1}b_{1}\ket{11}
    \end{aligned}
\end{equation}
This state $\ket{v_{A} v_{B}}$ is a linear combination of the four possible joint states of the two qubits. Each term in the expansion corresponds to one of the joint possibilities, and the coefficients $a_{0}b_{0}$, $a_{0}b_{1}$, $a_{1}b_{0}$, and $a_{1}b_{1}$ represent the probability amplitudes of each joint state.

\highspace
We are no longer describing two systems separately, we are describing \textbf{one system with 4 possible outcomes}.

\highspace
\begin{flushleft}
    \textcolor{Green3}{\faIcon{book} \textbf{General Two-Qubit State}}
\end{flushleft}
In \autoref{eq:two-qubit-state} we built a two-qubit state starting from a \textbf{very specific structure}, where the coefficients of the state were products of the coefficients of the individual qubits (i.e., $a_{i}b_{j}$). However, we can also consider a more \textbf{general two-qubit state}, where the coefficients are not necessarily products of the individual qubits' coefficients:
\begin{equation}\label{eq:general-two-qubit-state}
    \ket{v} = c_{0}\ket{00} + c_{1}\ket{01} + c_{2}\ket{10} + c_{3}\ket{11}
\end{equation}
Where $c_{0}$, $c_{1}$, $c_{2}$, and $c_{3}$ are \textbf{arbitrary complex numbers} that represent the probability amplitudes of each joint state.

\highspace
Obviously, this general two-qubit state \textbf{must satisfy the normalization condition}:
\begin{equation}
    \left|c_{0}\right|^{2} + \left|c_{1}\right|^{2} + \left|c_{2}\right|^{2} + \left|c_{3}\right|^{2} = 1
\end{equation}
This means that the sum of the probabilities of all joint states must equal 1, ensuring that the state is properly normalized.

\highspace
We can also \hl{generalize deeper to $n$ qubits}, where the state of the \textbf{system is a linear combination} of $2^n$ \textbf{basis states}, and the \textbf{coefficients} are arbitrary \textbf{complex numbers} that satisfy the normalization condition. Mathematically, we can express this as:
\begin{equation}\label{eq:general-multiple-qubit-state}
    \ket{v} = \sum_{k=0}^{2^n-1} c_k \ket{k}
\end{equation}
Where $\ket{k}$ represents the basis states of the $n$-qubit system, and $c_k$ are the corresponding probability amplitudes.  As we have already discussed at \autopageref{subsubsec:exponential-explosion}, this general form \hl{illustrates the exponential growth of the state space as the number of qubits increases}. However, we can write our one- and two-qubit states as special cases from this general form:
\begin{itemize}
    \item For $n=1$ qubit, we have $2^1 = 2$ basis states: $\ket{0}$ and $\ket{1}$, and the state can be written as:
    \begin{equation*}
        \ket{v} = c_0 \ket{0} + c_1 \ket{1}
    \end{equation*}
    \item For $n=2$ qubits, we have $2^2 = 4$ basis states: $\ket{00}$, $\ket{01}$, $\ket{10}$, and $\ket{11}$, and the state can be written as:
    \begin{equation*}
        \ket{v} = c_0 \ket{00} + c_1 \ket{01} + c_2 \ket{10} + c_3 \ket{11}
    \end{equation*}
\end{itemize}

\highspace
In summary, a general two-qubit state is a normalized superposition (i.e., the \emph{coefficients satisfy the normalization condition}) of the four basis states (i.e., $\ket{00}$, $\ket{01}$, $\ket{10}$, and $\ket{11}$). Unlike tensor product states, its coefficients are arbitrary and \emph{may not} factorize, leading to entanglement and non-classical correlations between the qubits (like the state in Example at \autopageref{ex:non-separable}).